\section{Introduction}

An agent that can accumulate useful assets must also decide when to spend, when to preserve liquidity, what to tell a counterparty and whether an apparent constraint is real. Over a long interaction, these decisions change the environment in which later decisions are made. An early bargain can fund a rival's construction; a profitable group can yield no paid rent before its owner is eliminated; a recovery policy can alter both the applied action and the memory available on the next turn. Evaluating the resulting agent requires more than checking whether its last action succeeded.

Long-horizon economic benchmarks make this problem concrete. Vending-Bench studies autonomous operation of a vending business, while its multiplayer Arena extension exposes agents to trade, strategic communication and competition~\cite{backlund2025vendingbench,li2026arena}. AgentMonopoly and FinBoardBench study economic choices in board-game environments, and DECE evaluates complete model--harness configurations under controlled economic conditions~\cite{wu2026agentmonopoly,hu2026finboardbench,zhang2026dece}. These studies already establish multiplayer economic interaction, legal-action interfaces, replay and state-checked communication as relevant evaluation tools. The remaining difficulty is connecting an observed outcome to the precise opportunities, information and resolved actions that produced it.

We study that connection in MonopolyBench, a reproducible four-player environment with persistent economic state, negotiation and an inspectable record of every decision. Our broad motivation is to evaluate both economic capability and alignment as agents become more capable. The present empirical question is narrower: \emph{what do complete competitive trajectories reveal about agents' acquisition, financing, development and social conduct, and which interpretations survive an audit of the alternatives and evidence available at the time?} A further design question asks what additional conditions would make the capability--alignment relationship identifiable.

We analyze twenty completed games containing 6,989 decisions and 8,679 response attempts. Four configured participants each play five times in every seat across five seed blocks. The record links applied actions and economic events to legal menus, delivered memory and available response channels. We combine event reconstruction, overlapping source review and exact-prefix execution of local alternatives. This permits a claim such as ``the debtor could pay'' to be tested as a legal sequence, and a claim such as ``the buyer was persuaded'' to be checked against prior motives, recipient exposure and the actual response.

The results support a detailed capability profile rather than a ranking. Trades often assemble groups, but two winners complete no trade. Development can be deferred with a stated reserve rationale, lost through response resolution, or inaccessible until financing changes. Three of sixty eliminations admit a legal route paying the current debt; another has substantial nominal value trapped by a hotel-liquidation rule. In the social record, misleading framing and selective disclosure coexist with explicit recipient correction, transient false echoes and bargaining that proceeds without identified uptake of the false premise. Since public deception is permitted, these observations describe conduct under a specific contract rather than a general instruction-misalignment rate.

Our contributions are threefold. First, we provide an integrated account of complete economic and social trajectories, preserving acquisition routes, bilateral transaction consequences, financing access, response provenance and downstream outcomes. Second, we apply an exact-prefix adjudication protocol to both agents' claims and generated research labels, distinguishing current menus from reachable legal sequences and recorded speech from delivered information and reliance. Third, we validate the recorded system and identify concrete requirements for future norm-explicit evaluation and training environments. The contribution is this combined protocol and empirical account, not a claim to have introduced replay, multiplayer economics or factual checking of agent communication.

